\subsection{DESTRUCTION}

The Shenandoah area, as a result of numerous campaigns and battles that took place there, suffered a great deal of damage. Some of this was the incidental damage of battle, but much other damage was a result of planned demolitions by both sides.

\begin{enumerate}
  \item Destruction can only be accomplished by Infantry, Cavalry and/or Partisan units.

  \item Destruction of unit counters such as Wagon or Supply counters is represented by removing them from the Map. Destruction to structures on the Mapboard (bridges and rail lines) is represented by placing a Destruction counter in the hex where the destruction occurred.

  \item Units can only carry out destruction in a hex they are in, never in an adjacent hex. A minimum of one Combat Factor is required for destruction.

  \item Destruction costs Movement Allowance Factors. (See the Destruction column of the MOVEMENT \& TERRAIN EFFECTS CHART.)

  \item Destruction cannot be carried out by a unit in Enemy Zone of Control, and can only be done by the Phasing Player.

  \item Wagon or Supply counters have to be destroyed individually. For example, one unit could expend Movement Allowance Factors once to destroy a Wagon counter, then again to destroy the Supply counter; or two units could both expend Movement Allowance Factors once each to perform one act of destruction each.

  \item Rail lines must be destroyed one hex at a time. If a rail hex also includes a hexside with a bridge, the bridge would have to be destroyed separately (place two Destruction Counters in the hex to signify that both bridge and and rail line are destroyed).

  \item A bridge can be destroyed from any hex that contains the hexside with the bridge.
\end{enumerate}

\subsection{CONSTRUCTION \& REPAIR}

\begin{enumerate}
  \item Construction and Repair can only be done by Infantry units.

  \item Destroyed Wagon and Supply counters cannot be repaired, they are out of the game for good. The repair of Mapboard structures is represented by removing the Destruction Counter.

  \item Units can only carry out Construction and Repair in a hex they occupy, never in an adjacent hex. A minimum of one Combat Factor is required for Construction and Repair.

  \item Construction and Repair costs Movement Allowance Factors (see the Construction \& Repair column of the MOVEMENT \& TERRAIN EFFECTS CHART).

  \item Construction and Repair cannot be carried out by a unit in an enemy Zone of Control, and can only be done by the Phasing Player.

  \item All Construction and Repair is performed individually on single structures. For instance, rail lines must be Repaired one hex at a time. If a rail hex also includes a bridge hexside, this would have to Repaired separately.

  \item To be Repaired, a bridge must have units in both hexes common to its hexside, and units in both of these hexes must expend Movement Allowance Factors to repair the one bridge.

  \item Fortifications can be constructed in a hex by expending Movement Allowance Factors for its Construction, then placing a Fortifications counter in the hex.

  \item No new bridges can be constructed, only old ones Repaired.

  \item Note on the MOVEMENT \& TERRAIN EFFECTS CHART the numbers in parenthesis under the Construction \& Repair column. This is the costs for repairing bridges only. All other Construction \& Repair uses the other costs.

  \item If any unit in a Stack is engaged in Construction \& Repair, all units in the Stack expend Movement Allowance Factors for the activity.
\end{enumerate}